\documentclass[11pt,a4paper]{article}

\usepackage[utf8]{inputenc} % allow utf-8 input
\usepackage[T1]{fontenc}    % use 8-bit T1 fonts
\usepackage{hyperref}       % hyperlinks
\usepackage{url}            % simple URL typesetting
\usepackage{booktabs}       % professional-quality tables
\usepackage{amsfonts}       % blackboard math symbols
\usepackage{nicefrac}       % compact symbols for 1/2, etc.
\usepackage{microtype}      % microtypography
\usepackage{xcolor}         % colors
\usepackage{graphicx}       % graphics
\usepackage[normalem]{ulem} % strikethrough
\title{Code snippets}

\begin{document}

\maketitle

The Pythagorean theorem can be written as an equation relating the lengths of the sides \textit{a} , \textit{b} and the hypotenuse \textit{c} .

To use Docling, simply install \texttt{docling} from your package manager, e.g. pip: \texttt{pip install docling}

To convert individual documents with python, use \texttt{convert()} , for example:

\begin{verbatim}
from docling.document_converter import DocumentConverter

source = "https://arxiv.org/pdf/2408.09869"
converter = DocumentConverter()
result = converter.convert(source)
print(result.document.export_to_markdown())
\end{verbatim}

The program will output: \texttt{\\#\\# Docling Technical Report[...]}

Prefetch the models:

\begin{itemize}
\item Use the \texttt{docling-tools models download} utility:
\item Alternatively, models can be programmatically downloaded using \texttt{docling.utils.model_downloader.download_models()} .
\item Also, you can use download-hf-repo parameter to download arbitrary models from HuggingFace by specifying repo id: \texttt{$ docling-tools models download-hf-repo ds4sd/SmolDocling-256M-preview Downloading ds4sd/SmolDocling-256M-preview model from HuggingFace...}
\end{itemize}

\end{document}